\problemname{Siffersumma}
\noindent
Siffrid älskar att leka med tal och har fått ett heltal $N$. 

Hon undrar hur man kan skapa det minsta och största talet som har samma siffersumma och antal siffror som $N$.

Siffersumman av ett tal $N$ är definierat som summan av alla siffror i talet. 

\noindent
Till exempel är siffersumman av talet $1234$: 
$$
1 + 2 + 3 + 4 = 10,
$$

och siffersumman av talet $220$:
$$
2 + 2 + 0 = 4.
$$

Kan du hjälpa Siffrid att hitta det minsta samt det största talet som har samma siffersumma och antal siffror som talet $N$?

\section*{Indata}
Indatan består av en rad med ett heltal, $N$ ($1 \le N \le 10^9$).

\section*{Utdata}
Skriv ut två heltal på samma rad, det minsta och det största talet med samma siffersumma och antal siffror som $N$.
Talen ska vara separerade med ett mellanslag.

\section*{Poängsättning}
Din lösning kommer att testas på en mängd testfallsgrupper.
För att få poäng för en grupp så måste du klara alla testfall i gruppen.

\noindent
\begin{tabular}{| l | l | p{12cm} |}
  \hline
  \textbf{Grupp} & \textbf{Poäng} & \textbf{Gränser} \\ \hline
  $1$    & $40$       & $N \leq 10^5$ \\ \hline
  $2$    & $60$       & Inga ytterligare begränsningar. \\ \hline
\end{tabular}

